\section{Validation Assets and Current Evidence}

\subsection{Validation Philosophy}

The current repository already contains a meaningful cross-validation suite, but
its evidence is distributed across individual test binaries rather than a single
benchmark report.  This notebook therefore documents the current validation
assets by test intent and by covered module interaction.

\subsection{Latest Release-Mode Core Snapshot}

After the recent power-flow and OPF refactors, the current release-mode evidence
set for the core notebook modules was rerun with the nine targeted binaries
listed below.  The run completed with 865 assertions and no failures.

\begin{longtable}{@{}L{0.32\linewidth}L{0.16\linewidth}L{0.42\linewidth}@{}}
\toprule
Test binary & Assertions & Main evidence added to this notebook \\
\midrule
\endhead
\texttt{test\_all\_components\_pf} & 159 & Full-component PF coverage after the \texttt{power\_flow} directory split and Jacobian-kernel refactor. \\
\texttt{test\_annual\_production\_sim} & 405 & Annual production simulation, L0/L2/L3 orchestration, and metric aggregation. \\
\texttt{test\_carbon\_analysis} & 82 & Static hybrid carbon tracing, matrix carbon intensity, and result summaries. \\
\texttt{test\_core\_modules} & 11 & Core module smoke and utility checks. \\
\texttt{test\_modeling\_toolkit} & 146 & AML/modeling-toolkit expression, bridge, and solver-facing contracts. \\
\texttt{test\_opf\_pf\_crossval} & 24 & OPF dispatch replay through PF feasibility checks. \\
\texttt{test\_powerflow\_performance} & 3 & Optimized PF kernel smoke/performance checks. \\
\texttt{test\_problem\_types} & 14 & Engine problem-type contract coverage. \\
\texttt{test\_projection\_equivalence} & 21 & Canonical projection equivalence and stable-result behavior. \\
\textbf{Total} & \textbf{865} & \textbf{All assertions passed.} \\
\bottomrule
\end{longtable}

\subsection{Cross-Validation Map}

\begin{longtable}{L{0.34\linewidth}L{0.54\linewidth}}
\toprule
Test or evidence source & What it validates \\
\midrule
\endhead
\texttt{tests/test\_opf\_pf\_crossval.cpp} & Consistency between AC OPF outputs and PF feasibility on representative AC/DC cases such as \texttt{ieee14\_acdc}, \texttt{case33bw\_acdc}, and \texttt{case33mg\_acdc}. \\
\texttt{tests/test\_ac\_opf\_jacobian\_fd.cpp} & Finite-difference verification of the native AC OPF equality Jacobian assembly. \\
\texttt{tests/test\_parity\_formulation\_fd.cpp} and \texttt{tests/test\_parity\_hessian\_fd.cpp} & Finite-difference verification of the parity formulation and its Hessian terms. \\
\texttt{tests/test\_parity\_ipm.cpp} and \texttt{tests/test\_parity\_julia.cpp} & Parity-IPM convergence and parity-path cross-checks against reference behavior. \\
\texttt{tests/test\_distribution\_pf\_crossval.cpp} & BFS-vs-Newton agreement on radial networks, plus validation of DER injections, hybrid feeder overlays, PV power curves, and stationary vs transit mobile storage behavior. \\
\texttt{tests/test\_short\_circuit\_crossval.cpp} & Analytical and self-consistency checks for short-circuit levels, Thevenin impedances, and API consistency between per-bus and batch runs. \\
\texttt{tests/test\_topology\_crossval.cpp} & Connectivity/radiality checks, 6-bus ONR comparison against exhaustive spanning-tree enumeration, and IEEE-33 style post-reconfiguration verification with Newton PF. \\
\texttt{tests/test\_solver\_handle.cpp} & Consistency between cached solver-handle runs and direct API runs. \\
\texttt{tests/test\_case\_builders.cpp}, \texttt{tests/test\_json\_roundtrip\_new\_components.cpp}, \texttt{tests/test\_excel\_comprehensive\_case.cpp} & Input-model integrity and serialization round-trip coverage for the expanded information model. \\
\texttt{tests/test\_acopf\_dcopf\_crossval.cpp} & DC OPF convergence on MATPOWER cases (IEEE 14/30/57/118/300), objective and dispatch cross-validation against AC OPF, LMP extraction, and PF feasibility check of DC OPF dispatch. \\
\texttt{tests/test\_annual\_production\_sim.cpp} & Annual production-simulation smoke tests: L0 annual planning, L2 weekly UC, L3 daily replay pipeline, and annual metric aggregation on small benchmark systems. \\
\texttt{tests/test\_lifecycle\_smoke.cpp} & Multi-year lifecycle simulation smoke tests: degradation application, stratified sampling, NPV accumulation, and year-by-year result consistency. \\
\bottomrule
\end{longtable}

\subsection{What Is Well Covered Today}

The current repository has comparatively strong evidence in the following areas.
\begin{itemize}
  \item Main PF correctness on representative AC/DC test cases.
  \item Derivative correctness for both native and parity OPF paths, including
  finite-difference verification of DCDC converter \(I^2R\) Jacobian and Hessian
  terms.
  \item VSC converter loss model consistency between PF and OPF
  (\(P_{\mathrm{loss}} = a + b\,I_{ac} + c\,I_{ac}^2\)).
  \item DCDC converter model consistency: efficiency, \(I^2R\) conduction loss,
  and correct sign-convention handling across PF and OPF solvers.
  \item OPF\(\to\)PF cross-validation on hybrid AC/DC cases, confirming
  that OPF dispatch is PF-feasible within documented tolerances.
  \item DC OPF vs AC OPF cross-validation on standard MATPOWER cases,
  confirming dispatch and objective agreement within expected DC-approximation
  error bounds, and PF feasibility of DC OPF dispatch.
  \item DPF agreement with Newton PF on radial benchmark feeders.
  \item ONR feasibility and post-solve verification through the full PF solver.
  \item Short-circuit self-consistency and small-network analytical checks.
  \item Input/output round trips for the enlarged public information model.
  \item Annual production simulation end-to-end pipeline (L0\(\to\)L2\(\to\)L3) on small benchmark systems.
  \item Lifecycle simulation multi-year degradation and NPV consistency.
\end{itemize}

\subsection{What Still Needs Better Validation}

The code structure itself highlights the most important evidence gaps.
\begin{itemize}
  \item DPF and ONR now share the canonical projection path, but the repository still lacks explicit projected-vs-canonical equivalence tests for feeders containing transformers, switches, flexible loads, asymmetric loads, and chargers.
  \item The detailed IEC short-circuit path is now exercised by dedicated tests, but stronger reference-case validation is still needed for explicit external-grid contribution accounting and richer component-by-component breakdown semantics.
  \item The parity path now instantiates richer component-variable blocks including DCDC converter dispatch with analytical Jacobian and Hessian, but the public OPF export contract still lags the in-memory result struct, so result-I/O regression tests should be extended when those fields are serialized.
  \item A direct notebook-generation regression target should eventually verify that all documented formulas still match the implementation.
  \item The reactive-power optimisation (RPO) path has no dedicated test file; convergence and OA gap properties are exercised only indirectly through internal development runs.
  \item The annual production simulation has smoke-level coverage but no reference-case comparison against an external production-cost tool.
  \item The lifecycle simulation has smoke-level coverage; stronger validation would include comparison of NPV results against manual spreadsheet calculations on a known degradation curve.
\end{itemize}

\subsection{Suggested Interpretation for Users}

At the current maturity level, the strongest production path is:
\begin{align}
\text{projected main PF / native OPF / parity OPF / DC OPF / batch+detailed SC / annual sim},
\end{align}
while DPF and ONR should still be used with awareness of their narrower hybrid
overlay coverage and their currently lighter result contracts.  This is not a
weakness of the numerical cores themselves; it is primarily an interface- and
consumer-contract issue between the rich public model layer and specialized
distribution-planning workflows.
